% Chapter 2 -- Literature Review (O-RAN / BMPP-DQN track)
%
% Governs: manuscript/ORAN_BMPP_DQN_Concept_Note_v1.md (Sections 3, 9)
% Chapter mapping: docs/oran_thesis_guide.md ("2. Literature Review" row)
% Status: DRAFT. Intended to be \input{} from the thesis-wide main.tex once
% that skeleton exists, immediately after chapter1_oran_introduction.tex
% (\label{chap:oran-introduction}). A thesis-wide references.bib does not
% yet exist, so citations are numbered locally (IEEE style) against the
% reference list at the end of this file, which repeats the keys already
% used in Chapter 1 plus the additional sources this chapter draws on.
% Merge into a single thesis-wide bibliography once main.tex is created.
%
% Scope note: parameter-validation literature (power-model constants,
% traffic-model breakpoints, functional-split bandwidth figures --
% manuscript/ORAN_BMPP_DQN_Concept_Note_v1.md Section 10's nine literature-
% check passes) is deliberately NOT reproduced here. Per
% docs/oran_thesis_guide.md's chapter-content mapping, that evidence belongs
% to Chapter 3 (System Model and Problem Formulation), where it supports
% specific numeric choices in oran_env/. This chapter instead evaluates the
% ALGORITHMIC and CONCEPTUAL literature that establishes the three research
% gaps introduced in Chapter 1.

\chapter{Literature Review}
\label{chap:oran-literature}

This chapter reviews the literature underlying the three research gaps
introduced in Chapter~\ref{chap:oran-introduction}: the absence of a
principled treatment of O-RAN's hybrid discrete-continuous action space
(Section~\ref{sec:lit-oran-energy}), the lack of an architecture that
combines action-branching decomposition with multi-pass parameterized
processing (Section~\ref{sec:lit-drl-hybrid}), and the limited treatment of
O-RAN's natural multi-timescale control structure
(Section~\ref{sec:lit-multitimescale}). Rather than surveying each paper in
isolation, the chapter is organized to build a cumulative argument: each
section evaluates what a cluster of related works does and does not
establish, so that Section~\ref{sec:lit-synthesis} can state precisely which
gap remains open after each cluster is accounted for.

\section{O-RAN Architecture and the Energy-Efficiency Imperative}
\label{sec:oran-arch-review}

The O-RAN Alliance's own architectural specification~\cite{oran2021mvp}
formalizes the RU/DU/CU disaggregation and the two-timescale RIC structure
introduced in Section~\ref{sec:oran-background}, and explicitly identifies
energy efficiency as a target use case for both RIC tiers. Subsequent survey
work has consistently reported the RAN as the dominant contributor to
mobile-network energy consumption and has catalogued AI/ML-driven control as
the primary proposed remedy: Barker, Seyfi, and Afghah's review of AI/ML in
mobile edge computing and Open RAN~\cite{barker2025mec} and Lassoued and
Boujnah's comprehensive review of 5G energy-efficiency
strategies~\cite{lassoued2026review} both position learning-based RAN
control as displacing static, rule-based energy-saving features (e.g.,
fixed sleep-mode schedules), on the grounds that traffic and channel
conditions are too non-stationary for a fixed policy to remain
near-optimal. Sthankiya et al.'s survey specifically of AI-driven energy
optimization in next-generation RAN~\cite{sthankiya2024survey} narrows this
further to O-RAN, and is useful for this thesis because it explicitly
catalogues which RIC tier, which disaggregation boundary, and which
learning paradigm (supervised, RL, federated) each surveyed work targets --
a taxonomy this review adopts implicitly in the sections that follow.
None of these three surveys, however, proposes or evaluates a specific
learning architecture; they establish that the problem is real and
important, and that RL is a leading candidate approach, but the specific
architectural question this thesis addresses -- how to represent O-RAN's
hybrid, multi-timescale action space inside an RL agent -- is left open by
design, as a survey's remit.

\section{Deep Reinforcement Learning for Discrete, Continuous, and Hybrid Action Spaces}
\label{sec:lit-drl-hybrid}

The DRL algorithms available to address O-RAN's control problem fall into
three families, and the boundary between them is the technical crux of this
thesis.

\textbf{Purely discrete-action methods.} Deep Q-Networks and their
double/dueling refinements learn a value function over a discrete action
set and select actions by maximization. In the adjacent C-RAN literature,
Iqbal, Tham, and Chang~\cite{iqbal2021ddqn} apply a Double DQN specifically
to RRH on/off decisions, pairing it with a separate convex solver for power
allocation -- a two-stage design this thesis's own discrete branches draw a
Double-DQN-style target computation from (Section~\ref{sec:oran-background}
already anticipates this; the mechanism is specified formally in
Chapter~\ref{chap:oran-system-model}). Fathy, Abood, and
Hamdi~\cite{fathy2021ann} similarly pair a neural network (there, used for
pre-processing rather than value estimation) with a convex power-allocation
stage. Both works demonstrate that a discrete activation decision and a
continuous power decision can be decoupled into two separately-solved
stages, but neither trains the two stages jointly, end-to-end, inside a
single RL agent -- the coupling between them is handled procedurally
(the discrete stage's output is simply passed as an input to the convex
stage), not learned. This two-stage decoupling is a reasonable engineering
simplification when a convex solver is available for the continuous
sub-problem, but O-RAN's own energy-power relationship (transmit power
interacting with PRB share, interference, and functional-split-dependent
processing cost) is not naturally convex in the way idealized power-control
problems are, which is part of the motivation for a jointly-learned
approach.

\textbf{Purely continuous-action methods.} Deep Deterministic Policy
Gradient (DDPG) and its refinements learn a deterministic policy over a
continuous action space via a differentiable critic. These methods have no
native mechanism for a discrete choice: a binary decision such as RU
activation can only be approximated by thresholding a continuous output,
which -- as already noted in Section~\ref{sec:oran-problem-statement} for
OREO~\cite{qazzaz2026oreo} specifically -- breaks the gradient signal a true
discrete decision would carry, since the threshold operation's gradient is
zero almost everywhere. Shengren et al.'s benchmark of DRL algorithms for
energy-systems scheduling~\cite{shengren2022benchmark} (SAC $>$ TD3 $>$
DDPG in training stability, evaluated on a purely continuous control
problem) is useful evidence for algorithm selection \emph{within} the
continuous family, but by construction says nothing about the
discrete-continuous coupling problem, since its own benchmark task has no
discrete component.

\textbf{Parameterized (hybrid) action-space methods.} P-DQN, introduced by
Xiong et al.~\cite{xiong2018pdqn}, is the first architecture designed
specifically for actions that are simultaneously discrete (which action to
take) and continuous (a parameter of that action): a shared network
produces continuous parameters for every discrete action in parallel, and a
Q-network scores each (discrete action, continuous parameter) pair. This
directly targets the coupling problem left unaddressed by the two clusters
above. However, P-DQN's own single-pass design has a documented flaw: the
Q-network's input includes \emph{every} discrete action's continuous
parameters simultaneously, so the gradient used to train the parameter
network for action $k$ is contaminated by the (irrelevant) parameters of
every other action $k' \neq k$ -- a false cross-talk gradient rather than a
genuine learning signal. Bester, James, and Konidaris's MP-DQN
refinement~\cite{bester2019mpdqn} fixes exactly this flaw with multi-pass
processing: each discrete action's Q-value is computed from a masked input
in which only that action's own continuous parameters are non-zero,
eliminating the cross-talk gradient by construction. This progression
(P-DQN $\to$ MP-DQN) is the direct methodological ancestor of this thesis's
own multi-pass masking mechanism (Chapter~\ref{chap:oran-system-model}).

Two applied parameterized-action papers confirm this family's relevance to
wireless resource control specifically, while also exposing its remaining
limitation. Sharma and Yoon~\cite{sharma2023massivemimo} apply a
parameterized DQN to power allocation in massive MIMO, and Kabir, Tham, and
Chang~\cite{kabir2023mobility} apply parameterized RL to mobility-aware IoT
resource allocation; a broader, unrelated IEEE study of hybrid-action RL
for load balancing and energy efficiency~\cite{ieee2024selfoptimized} and
Yan et al.'s constrained hybrid-action-space
formulation~\cite{yan2025fluctuates} further confirm that parameterized RL
is an active, validated methodology beyond this thesis's specific setting.
None of these four works, however, uses a \emph{branching} decomposition of
the discrete side of the action space: each treats the discrete action set
as a single flat choice (e.g., ``which sub-channel,'' ``which mobility
mode''), not as a vector of $R$ independent per-entity decisions the way RU
activation naturally is (one binary decision per RU). This is precisely
where P-DQN and MP-DQN's remaining scalability limitation becomes material
for O-RAN: with $R$ RUs, $S$ split options, and a flat parameterized
representation, the discrete action set already grows combinatorially
(e.g., $S^R$ joint split choices before RU activation is even considered),
reintroducing at the architectural level the same combinatorial blow-up
that motivated moving from tabular Q-learning to function approximation in
the first place.

Branching DQN (BDQ), proposed by Tavakoli, Pardo, and
Kormushev~\cite{tavakoli2018branching}, addresses exactly this combinatorial
growth, but for purely discrete action spaces: by giving each of $R$
independent decisions its own ``branch'' (a dueling-style Q-head sharing a
common state encoder), BDQ reduces the discrete output size from
$\prod_r |\mathcal{A}_r|$ to $\sum_r |\mathcal{A}_r|$ -- linear rather than
exponential in $R$. This is the second direct methodological ancestor of
this thesis's design. Critically, however, BDQ's own formulation has no
continuous-parameter mechanism at all: it is a discrete-only architecture,
and its branches' independence (each branch's Q-value depends only on the
shared state encoding, not on any other branch's chosen action) is a
design choice made for a purely discrete setting, where it is a reasonable
one. Chapter~\ref{chap:oran-results} returns to this specific independence
assumption when discussing what the empirical comparison in this thesis
reveals about its consequences in an O-RAN environment where cross-RU
interference genuinely couples the branches' outcomes.

\section{Deep Reinforcement Learning for O-RAN Energy Optimization}
\label{sec:lit-oran-energy}

Having established the three architectural families in isolation, this
section evaluates how each has actually been applied to O-RAN (or
sufficiently O-RAN-adjacent) energy optimization specifically, and where
each application still leaves the discrete-continuous, multi-branch,
multi-timescale coupling problem unresolved.

OREO~\cite{qazzaz2026oreo} remains the most directly comparable prior
system: it targets the same problem (O-RAN energy optimization via DRL) at
the same architectural ambition (a learned policy, not a heuristic), but --
as detailed in Sections~\ref{sec:oran-problem-statement} and
\ref{sec:lit-drl-hybrid} -- resolves the hybrid-action problem only by
approximation (continuous relaxation plus thresholding), uses a flat
(non-branching) action representation, and operates its PPO policy at a
single timescale (the Non-RT RIC's), leaving Near-RT-RIC-scale decisions and
any coupling between the two timescales outside its scope. Amiri, Wang,
Shojafar, and Tafazolli~\cite{amiri2022vnf} apply DRL to dynamic VNF
(functional-split-related) splitting decisions in O-RAN specifically, which
is methodologically close to this thesis's own split-selection branch, but
their formulation treats split selection as the sole learned decision and
does not jointly optimize it alongside RU activation or continuous power/PRB
control, so the cross-decision coupling this thesis targets is outside
their scope by construction, not merely unaddressed. Bordin et
al.~\cite{bordin2025design} evaluate both PPO and DQN for RF-front-end
energy saving in Open RAN using the \texttt{ns-O-RAN} simulator, providing
useful evidence that both algorithm families are viable in a realistic
O-RAN simulation environment, but their formulation, like OREO's, does not
combine branching decomposition with parameterized continuous control; it
compares algorithm families on a task that itself remains discrete
(RF-front-end on/off), so it does not test the hybrid-action question this
thesis addresses. Rezazadeh, Chergui, Christofi, and Verikoukis's
actor-critic approach to zero-touch resource and energy control in network
slicing~\cite{rezazadeh2022actorcritic} confirms actor-critic methods are
viable for joint resource-energy control in a RAN-adjacent setting, but
targets slice-level resource control rather than O-RAN's specific
RU/DU/CU-disaggregated action space, and does not decompose its action
space by branch at all.

The closest related work to this thesis's own proposed architecture is Lu,
Yan, and Zeng's EExApp~\cite{luyanzeng2026eexapp}, which addresses RU
energy optimization in 5G O-RAN using a GNN-based RL formulation with two
separate actor-critic pairs -- one for RU sleep decisions, one for DU
network slicing -- coordinated as two interacting agents rather than one.
EExApp is evidence that decomposing an O-RAN energy-control problem by
decision type (sleep vs.\ slicing) is a productive design choice, which
this thesis's own branch decomposition shares in spirit. The architectural
difference is nonetheless material: EExApp's two decision types are each handled by an independent actor-critic
pair with its own critic, whereas this thesis's BMPP-DQN handles both
decision types through a single \emph{critic} (a non-twin
\texttt{BranchingCritic}, evaluated once per branch via multi-pass masking)
rather than two separately-trained critics. BMPP-DQN does still use two
separate encoders internally -- an upper encoder feeding the discrete
branching critic and a lower encoder feeding the continuous parameter
network, reflecting the two-timescale design RQ3 investigates -- so the
distinction from EExApp is specifically in the \emph{value function}: one
shared critic scoring both decision types' branches, rather than two
independently-trained critics each scoring one decision type. This matters
for the coordination question this thesis targets: two independently
trained critics must coordinate through whatever channel connects their
respective agents (e.g., one agent's action entering the other's observed
state), whereas a single shared critic scores every branch's action against
the same learned value function directly. Whether a shared-critic design in
fact resolves EExApp's own two-agent coordination problem better than two
coordinated critics do is an empirical question this thesis's own results,
not this literature review, can speak to; the two designs represent
genuinely different resolutions of the same underlying decomposition idea,
not a strict improvement of one over the other on paper alone.

A further comparison is instructive for what it \emph{cannot} establish.
Chuang et al.'s hybrid A3C and Dueling DQN architecture~\cite{chuang2025hybrid}
combines two DRL paradigms in a 5G C-RAN setting, superficially resembling
this thesis's own combination of two DRL paradigms (branching, multi-pass
parameterized processing). On inspection, however, Chuang et al.\ address
industrial IoT task scheduling, not RU/DU/CU energy control, and their
hybrid combines two purely discrete-action methods (A3C's policy-gradient
updates and Dueling DQN's value estimation), not a discrete-continuous
hybrid action space. The paper is therefore evidence only that combining
two DRL paradigms within one architecture is a viable general strategy, a
premise this thesis also relies on; it does not bear on the specific
discrete-continuous, branch-structured coupling problem that motivates
BMPP-DQN, and is engaged here only to disclose that it was considered and
found not directly applicable, rather than silently omitted.

\section{Multi-Timescale Control in O-RAN}
\label{sec:lit-multitimescale}

The O-RAN Alliance's own two-tier RIC specification~\cite{oran2021mvp}
establishes that O-RAN's control problem is natively multi-timescale by
design, not merely as an implementation detail: the Non-RT and Near-RT RICs
are specified with different control-loop periods precisely because
different categories of decision (policy/configuration versus
resource/scheduling) have different natural update rates. Two further works
engage this structure without resolving it in the way this thesis's RQ3
requires. Liang, Al-Tahmeesschi, Chetty, and Ahmadi's federated-TD3
approach~\cite{liang2026federated} coordinates a rApp aggregator with
per-cell xApp agents -- a genuine multi-tier architecture -- but its
tiering separates \emph{agents by network entity} (a central aggregator
versus distributed per-cell learners) rather than \emph{decisions by
timescale} within a single agent; its motivating rationale (avoiding the
scalability cost of a single centralized agent as the number of cells
grows) is itself cited in Section~\ref{sec:oran-scope}'s justification for
this thesis's own single-agent scope, but the paper does not itself study
how separating decision timescales within one agent affects convergence.
Sohaib et al.'s transfer-learning approach for cloud-native
O-RAN~\cite{sohaib2024transfer} demonstrates that a policy trained under one
traffic condition (e.g., eMBB) can be adapted to another (e.g., URLLC) with
reduced retraining cost, which is a form of temporal adaptation, but again
addresses adaptation \emph{across} training episodes, not the
within-episode coupling between a slow branch's decision and a fast
branch's subsequent, dependent decisions that this thesis's two-timescale
design (RU activation and split selection held fixed for several
environment steps while power and PRB allocation update every step) is
built to capture. No work reviewed in this chapter evaluates, within a
single learned agent, how explicitly separating slow structural decisions
from fast resource decisions affects convergence relative to treating both
at a single update frequency -- the specific comparison this thesis's RQ3
poses.

\section{Synthesis: Research Gaps Addressed by This Thesis}
\label{sec:lit-synthesis}

Table~\ref{tab:lit-gap-synthesis} summarizes how the clusters reviewed above
map onto the three gaps introduced in Section~\ref{sec:oran-problem-statement}.

\begin{table}[htbp]
  \centering
  \caption{Mapping from reviewed literature clusters to the three research
    gaps this thesis addresses. A checkmark indicates the cluster resolves
    that dimension of the gap; a dash indicates it does not.}
  \label{tab:lit-gap-synthesis}
  \begin{tabular}{lccc}
    \toprule
    Cluster & Discrete-continuous & Branching & Multi-timescale \\
            & coupling (Gap 1) & decomposition (Gap 2) & coordination (Gap 3) \\
    \midrule
    DQN family~\cite{iqbal2021ddqn,fathy2021ann} & -- (two-stage, not joint) & -- & -- \\
    DDPG family~\cite{shengren2022benchmark} & -- (thresholding only) & -- & -- \\
    P-DQN / MP-DQN~\cite{xiong2018pdqn,bester2019mpdqn} & \checkmark & -- (flat) & -- \\
    Branching DQN~\cite{tavakoli2018branching} & -- (discrete only) & \checkmark & -- \\
    OREO~\cite{qazzaz2026oreo} & -- (thresholding) & -- (flat) & -- (single tier) \\
    EExApp~\cite{luyanzeng2026eexapp} & \checkmark (per agent) & \checkmark (2 agents) & -- \\
    Federated TD3~\cite{liang2026federated} & -- & -- & -- (agent-tiered, not decision-tiered) \\
    \textbf{BMPP-DQN (this thesis)} & \checkmark & \checkmark & \checkmark \\
    \bottomrule
  \end{tabular}
\end{table}

Three conclusions follow. First, the discrete-continuous coupling problem
(Gap 1) is solved in principle by the P-DQN/MP-DQN line, but no O-RAN energy
work reviewed here adopts it; O-RAN energy papers instead resolve the
hybrid action space either by thresholding a continuous policy (OREO) or by
decoupling discrete and continuous decisions into separately-trained
components (Iqbal et al., Fathy et al., in the adjacent C-RAN literature).
Second, the branching decomposition that would let a parameterized method
scale to $R$ independent RUs (Gap 2) exists for discrete-only action spaces
(BDQ) but has not been combined with multi-pass parameterized processing in
any work reviewed here -- EExApp's two-agent decomposition is the closest
attempt, but decomposes by decision \emph{type} across two independently
trained agents, not by \emph{entity} within one coupled network the way BDQ
decomposes by RU. Third, no work reviewed here evaluates, within a single
agent, the effect of explicitly separating decision timescales on
convergence (Gap 3); the closest analogues either tier by network entity
(federated TD3) or by training phase (transfer learning), not by decision
frequency within one training loop.

BMPP-DQN, introduced formally in Chapter~\ref{chap:oran-system-model},
is positioned directly at the intersection these three gaps define: it
extends BDQ's per-RU branching with MP-DQN's multi-pass masking so that each
branch's discrete Q-value is computed from a shared encoding without
cross-talk from other branches' continuous parameters, and it separates the
resulting branches into an upper-level (RU activation, split selection) and
lower-level (power, PRB allocation) network reflecting the Non-RT/Near-RT
RIC timescale split the O-RAN Alliance's own architecture
specifies~\cite{oran2021mvp}. Chapter~\ref{chap:oran-results} evaluates
this design against DQN, DDPG, and MP-DQN baselines and reports what the
resulting comparison actually shows about each of these three design
choices -- including, as Chapter~\ref{chap:oran-conclusion} discusses in
full, a finding about branch independence that this literature review's own
treatment of BDQ's discrete-only design in
Section~\ref{sec:lit-drl-hybrid} anticipates as a live architectural
question rather than a settled one.

% ---------------------------------------------------------------------------
% Local reference list for this chapter (IEEE numbered style). Entries
% correspond to manuscript/ORAN_BMPP_DQN_Concept_Note_v1.md Section 9 and
% AGENTS.md's Foundational References table. Keys already used in
% chapter1_oran_introduction.tex are repeated verbatim here so this chapter
% remains standalone-compilable; merge into a single thesis-wide
% references.bib once main.tex exists, and do not let the two chapters'
% copies of a shared key drift apart until then.
% ---------------------------------------------------------------------------
\begin{thebibliography}{21}

\bibitem{qazzaz2026oreo}
M.~M.~H. Qazzaz, A. Salama, M. Hafeez, and S.~A.~R. Zaidi,
``OREO: Open RAN Energy Optimization via Deep Reinforcement Learning for 6G
Networks,'' \emph{IEEE Open Journal of the Communications Society}, vol.~7,
pp.~4165--4182, 2026.

\bibitem{oran2021mvp}
O-RAN Alliance, ``O-RAN Minimum Viable Plan and Commercialization
Strategy,'' Jun. 2021.

\bibitem{xiong2018pdqn}
J. Xiong, Q. Wang, Z. Yang, P. Sun, L. Han, Y. Zheng, H. Fu, T. Zhang, J. Liu,
and H. Liu, ``Parametrized Deep Q-Networks Learning: Reinforcement Learning
with Discrete-Continuous Hybrid Action Space,'' arXiv preprint
arXiv:1810.06394, 2018.

\bibitem{bester2019mpdqn}
C.~J. Bester, S.~D. James, and G.~D. Konidaris, ``Multi-Pass Q-Networks for
Deep Reinforcement Learning with Parameterised Action Spaces,'' arXiv
preprint arXiv:1905.04388, 2019.

\bibitem{tavakoli2018branching}
A. Tavakoli, F. Pardo, and P. Kormushev, ``Action Branching Architectures for
Deep Reinforcement Learning,'' in \emph{Proc. AAAI Conference on Artificial
Intelligence}, 2018.

\bibitem{bordin2025design}
M. Bordin, A. Lacava, M. Polese, S. Satish, M. AnanthaSwamy Nittoor,
R. Sivaraj, F. Cuomo, and T. Melodia, ``Design and Evaluation of Deep
Reinforcement Learning for Energy Saving in Open RAN,'' in \emph{Proc. IEEE
22nd Consumer Communications \& Networking Conference (CCNC)}, Las Vegas,
NV, USA, Jan. 2025, pp.~1--6.

\bibitem{sthankiya2024survey}
A. Sthankiya et al., ``A Survey on AI-Driven Energy Optimization in
Next-Generation RAN,'' arXiv preprint arXiv:2411.02164, 2024.

\bibitem{liang2026federated}
X. Liang, A. Al-Tahmeesschi, S.~B. Chetty, and H. Ahmadi, ``Green O-RAN
Operation: a Modern ML-Driven Network Energy Consumption Optimization,'' in
\emph{Proc. IEEE GLOBECOM}, 2025, pp.~4475--4480.

\bibitem{sohaib2024transfer}
M. Sohaib et al., ``DRL Transfer Learning for Cloud-Native O-RAN,''
arXiv preprint arXiv:2407.11563, 2024.

\bibitem{iqbal2021ddqn}
A. Iqbal, M.-L. Tham, and Y.~C. Chang, ``Double Deep Q-Network-Based
Energy-Efficient Resource Allocation in Cloud Radio Access Network,''
\emph{IEEE Access}, vol.~9, pp.~20440--[end page not independently
confirmed in this repository], 2021, doi: 10.1109/ACCESS.2021.3054909.

\bibitem{fathy2021ann}
M. Fathy, M.~S. Abood, and M.~M. Hamdi, ``Optimization of Energy-Efficient
Cloud Radio Access Networks for 5G using Neural Networks,'' in \emph{Proc.
Int. Conf. Intelligent Technology, System and Service for Internet of
Everything (ITSS-IoE)}, 2021.

\bibitem{alzubaedi2019cran}
W.~H.~A. Al-Zubaedi, ``Planning a C-RAN Deployment for the Next Generation
Cellular Networks,'' Ph.D. dissertation, Dept.\ of Electronic and Computer
Engineering, Brunel University London, London, U.K., 2019.

\bibitem{shengren2022benchmark}
H. Shengren, E.~M. Salazar Duque, P.~P. Vergara, and P. Palensky,
``Performance Comparison of Deep RL Algorithms for Energy Systems Optimal
Scheduling,'' in \emph{Proc. IEEE PES Innovative Smart Grid Technologies
Conference Europe (ISGT-Europe)}, 2022, pp.~1--6.

\bibitem{ieee2024selfoptimized}
``Self-Optimized Agent for Load Balancing and Energy Efficiency: A
Reinforcement Learning Framework With Hybrid Action Space,'' \emph{IEEE
Trans.\ on Network and Service Management}, vol.~21, no.~4, pp.~4902--4919,
Jul. 2024.

\bibitem{sharma2023massivemimo}
S. Sharma and W. Yoon, ``Energy Efficient Power Allocation in Massive MIMO
Based on Parameterized Deep DQN,'' \emph{Electronics}, vol.~12, no.~21,
p.~4517, 2023.

\bibitem{kabir2023mobility}
H. Kabir, M.-L. Tham, and Y.~C. Chang, ``Mobility-Aware Resource Allocation
in IoT Network for Post-Disaster Communications with Parameterized
Reinforcement Learning,'' \emph{Sensors}, vol.~23, no.~14, p.~6448,
Jul. 2023.

\bibitem{luyanzeng2026eexapp}
J. Lu, P. Yan, and H. Zeng, ``EExApp: GNN-Based Reinforcement Learning for
Radio Unit Energy Optimization in 5G O-RAN,'' in \emph{Proc.\ IEEE
INFOCOM}, 2026, pp.~1--10.

\bibitem{yan2025fluctuates}
C. Yan et al., ``Hybrid Reinforcement Learning in Parameterized Action
Space via Fluctuates Constraint,'' \emph{Engineering Applications of
Artificial Intelligence}, vol.~162, p.~112499, 2025.

\bibitem{rezazadeh2022actorcritic}
F. Rezazadeh, H. Chergui, L. Christofi, and C. Verikoukis, ``Actor-Critic-
Based Learning for Zero-Touch Joint Resource and Energy Control in Network
Slicing,'' arXiv preprint, 2022.

\bibitem{amiri2022vnf}
E. Amiri, N. Wang, M. Shojafar, and R. Tafazolli, ``Energy-Aware Dynamic
VNF Splitting in O-RAN Using Deep Reinforcement Learning,'' in \emph{Proc.
IEEE Int.\ Conf.\ Local Computer Networks (LCN)}, 2022, pp.~422--429.

\bibitem{barker2025mec}
R. Barker, T. Seyfi, and F. Afghah, ``Advancements in Mobile Edge Computing
and Open RAN: Leveraging Artificial Intelligence and Machine Learning for
Wireless Systems,'' \emph{arXivLens}, 2025.

\bibitem{lassoued2026review}
N. Lassoued and N. Boujnah, ``A Comprehensive Review of Energy Efficiency in
5G Networks: Past Strategies, Present Advances, and Future Research
Directions,'' \emph{Computers}, vol.~15, no.~1, p.~50, 2026.

\bibitem{chuang2025hybrid}
Chuang et al., ``[Title and venue not independently verified in this
repository -- known only via a descriptive secondary reference, per
AGENTS.md's Foundational References table:] Hybrid A3C and Dueling DQN
applied to industrial IoT task scheduling in a 5G Cloud-RAN setting,''
2025. \emph{Needs primary-source verification before this entry is used
outside this draft.}

\end{thebibliography}
